\chapter{2002年大纲卷（理）}
\section{单选题}
\begin{problem}
圆 \(\left(x - 1\right)^2 + y^2 = 1\) 的圆心到直线 \(y = \frac{\sqrt{3}}{3} x\) 的距离是（\quad）

\choices
    {\( \frac{1}{2} \)}
    {\(\frac{\sqrt{3}}{2}\)}
    {\(1\)}
    {\(\sqrt{3}\)}
\end{problem}

\begin{answer}
A
\end{answer}

\begin{solution}
圆心为 \((1,0)\)，直线可化为 \(\sqrt{3}x-3y=0\). 因此圆心到该直线的距离为
\(\displaystyle\frac{|\sqrt{3}|}{\sqrt{(\sqrt{3})^2+(-3)^2}}=\frac{1}{2}\). 故选 A.
\end{solution}

\begin{problem}
复数 \(\left(\frac{1}{2} + \frac{\sqrt{3}}{2}\i\right)^3\) 的值是（\quad）

\choices
    {\(-\i\)}
    {\(\i\)}
    {\(-1\)}
    {\(1\)}
\end{problem}

\begin{answer}
C
\end{answer}

\begin{solution}
因为 \(\displaystyle\frac{1}{2}+\frac{\sqrt{3}}{2}\i=\cos\frac{\pi}{3}+\i\sin\frac{\pi}{3}\)，所以由棣莫弗定理，
\(\displaystyle\left(\frac{1}{2}+\frac{\sqrt{3}}{2}\i\right)^3=\cos\pi+\i\sin\pi=-1\). 故选 C.
\end{solution}

\begin{problem}
不等式 \(\left(1 + x\right)\left(1 - |x|\right) > 0\) 的解集是（\quad）

\choices
    {\(\{x\mid 0\leqslant x < 1\}\)}
    {\(\{x\mid x < 0 \text{ 且 } x\neq -1\}\)}
    {\(\{x\mid -1 < x < 1\}\)}
    {\(\{x\mid x < 1 \text{ 且 } x\neq -1\}\)}
\end{problem}

\begin{answer}
D
\end{answer}

\begin{solution}
当 \(x\geqslant0\) 时，原不等式化为 \((1+x)(1-x)>0\)，故 \(0\leqslant x<1\). 当 \(x<0\) 时，\(|x|=-x\)，原不等式化为 \((1+x)^2>0\)，故 \(x<0\) 且 \(x\neq-1\). 合并得解集为 \(\{x\mid x<1\text{ 且 }x\neq-1\}\)，故选 D.
\end{solution}

\begin{problem}
在 \((0,2\pi)\) 内，使 \(\sin x > \cos x\) 成立的 \(x\) 的取值范围是（\quad）

\choices
    {\(\left(\frac{\pi}{4},\frac{\pi}{2}\right)\cup \left(\pi,\frac{5\pi}{4}\right)\)}
    {\(\left(\frac{\pi}{4},\pi\right)\)}
    {\(\left(\frac{\pi}{4},\frac{5\pi}{4}\right)\)}
    {\(\left(\frac{\pi}{4}, \pi\right) \cup \left(\frac{5\pi}{4}, \frac{3\pi}{2}\right)\)}
\end{problem}

\begin{answer}
C
\end{answer}

\begin{solution}
由 \(\sin x-\cos x=\sqrt{2}\sin\left(x-\frac{\pi}{4}\right)\)，原不等式等价于 \(\sin\left(x-\frac{\pi}{4}\right)>0\). 在 \(0<x<2\pi\) 内，正弦值为正的部分为 \(\frac{\pi}{4}<x<\frac{5\pi}{4}\)，故选 C.
\end{solution}

\begin{problem}
设集合 \(M = \left\{x \mid x = \frac{k}{2} + \frac{1}{4}, k \in \Z\right\}\)，\(N = \left\{x \mid x = \frac{k}{4} + \frac{1}{2}, k \in \Z\right\}\)，则（\quad）

\choices
    {\(M = N\)}
    {\(M \subseteq N\)}
    {\(M \supseteq N\)}
    {\(M \cap N = \varnothing\)}
\end{problem}

\begin{answer}
B
\end{answer}

\begin{solution}
集合 \(M\) 中的元素可写成 \(\displaystyle\frac{2k+1}{4}\)，而集合 \(N\) 中的元素可写成 \(\displaystyle\frac{k+2}{4}\). 对任意 \(k\in\Z\)，取 \(l=2k-1\in\Z\)，则 \(\displaystyle\frac{2k+1}{4}=\frac{l+2}{4}\in N\)，所以 \(M\subseteq N\). 例如 \(0=\frac{-2+2}{4}\in N\) 但 \(0\notin M\)，故不是相等关系，选 B.
\end{solution}

\begin{problem}
点 \(P(1,0)\) 到曲线 \(\begin{cases}
x = t^2,\\
y = 2t,
\end{cases}\) （其中参数 \(t \in \R\)）上的点的最短距离为（\quad）

\choices
    {\(0\)}
    {\(1\)}
    {\(\sqrt{2}\)}
    {\(2\)}
\end{problem}

\begin{answer}
B
\end{answer}

\begin{solution}
曲线上点的坐标为 \((t^2,2t)\)，所以它到 \(P(1,0)\) 的距离的平方为
\(\displaystyle(t^2-1)^2+(2t)^2=t^4+2t^2+1=(t^2+1)^2\). 因为 \(t^2+1>0\)，距离为 \(t^2+1\geqslant1\)，且在 \(t=0\) 时取到，故最短距离为 \(1\)，选 B.
\end{solution}

\begin{problem}
一个圆锥和一个半球有公共底面，如果圆锥的体积恰好与半球的体积相等，那么这个圆锥轴截面顶角的余弦值是（\quad）

\choices
    {\(\frac{3}{4}\)}
    {\(\frac{4}{5}\)}
    {\(\frac{3}{5}\)}
    {\(-\frac{3}{5}\)}
\end{problem}

\begin{answer}
C
\end{answer}

\begin{solution}
设圆锥和半球的公共底面半径为 \(r\)，圆锥高为 \(h\). 由体积相等，
\(\displaystyle\frac{1}{3}\pi r^2h=\frac{2}{3}\pi r^3\)，得 \(h=2r\). 轴截面为腰长 \(\sqrt{r^2+h^2}=\sqrt{5}r\)、底边长 \(2r\) 的等腰三角形. 设其顶角为 \(\theta\)，由余弦定理，
\(\displaystyle\cos\theta=\frac{(\sqrt{5}r)^2+(\sqrt{5}r)^2-(2r)^2}{2(\sqrt{5}r)^2}=\frac{3}{5}\). 故选 C.
\end{solution}

\begin{problem}
正六棱柱 \(ABCDEF - A_{1}B_{1}C_{1}D_{1}E_{1}F_{1}\) 的底面边长为 \(1\)，侧棱长为 \(\sqrt{2}\)，则这个棱柱侧面对角线 \(E_{1}D\) 与 \(BC_{1}\) 所成的角是（\quad）

\choices
    {\(90^{\circ}\)}
    {\(60^{\circ}\)}
    {\(45^{\circ}\)}
    {\(30^{\circ}\)}
\end{problem}

\begin{answer}
B
\end{answer}

\begin{solution}
建立直角坐标系，使正六边形底面位于 \(z=0\) 平面，并取
\(\displaystyle A=(1,0,0)\)，\(\ B=\left(\frac12,\frac{\sqrt3}{2},0\right)\)，\(\ C=\left(-\frac12,\frac{\sqrt3}{2},0\right)\)，\(\ D=(-1,0,0)\)，\(\ E=\left(-\frac12,-\frac{\sqrt3}{2},0\right)\). 侧棱方向为 \((0,0,\sqrt2)\)，于是
\(\displaystyle\overrightarrow{DE_1}=\left(\frac12,-\frac{\sqrt3}{2},\sqrt2\right)\)，
\(\displaystyle\overrightarrow{BC_1}=\left(-1,0,\sqrt2\right)\). 因而
\(\displaystyle\overrightarrow{DE_1}\cdot\overrightarrow{BC_1}=\frac32\)，且 \(\left|\overrightarrow{DE_1}\right|=\left|\overrightarrow{BC_1}\right|=\sqrt3\). 若所成的角为 \(\theta\)，则 \(\cos\theta=\frac{1}{2}\)，故 \(\theta=60^{\circ}\)，选 B.
\end{solution}

\begin{problem}
函数 \(y = x^{2} + bx + c\)（\(x \in [0, +\infty)\)）是单调函数的充要条件是（\quad）

\choices
    {\(b \geqslant 0\)}
    {\(b \leqslant 0\)}
    {\(b > 0\)}
    {\(b < 0\)}
\end{problem}

\begin{answer}
A
\end{answer}

\begin{solution}
函数在 \([0,+\infty)\) 上的导数为 \(y'=2x+b\). 若 \(b\geqslant0\)，则对所有 \(x\geqslant0\) 有 \(y'\geqslant0\)，函数在该区间上单调递增. 若 \(b<0\)，令 \(x_0=-\frac b2>0\)，则当 \(0\leqslant x<x_0\) 时有 \(y'<0\)，当 \(x>x_0\) 时有 \(y'>0\)，函数先减后增，因而既不是单调递增函数，也不是单调递减函数. 所以充要条件为 \(b\geqslant0\)，选 A.
\end{solution}

\begin{problem}
函数 \( y = 1 - \frac{1}{x - 1} \) 的图象是（\quad）

\choices
    {\choicebitmap[width=0.15\paperwidth]{img/2002/syllabus_science/q10_optA.png}}
    {\choicebitmap[width=0.15\paperwidth]{img/2002/syllabus_science/q10_optB.png}}
    {\choicebitmap[width=0.15\paperwidth]{img/2002/syllabus_science/q10_optC.png}}
    {\choicebitmap[width=0.15\paperwidth]{img/2002/syllabus_science/q10_optD.png}}
\end{problem}

\begin{answer}
B
\end{answer}

\begin{solution}
函数满足 \(\displaystyle(x-1)(y-1)=-1\)，所以图象的两条渐近线为 \(x=1\) 和 \(y=1\). 当 \(x<1\) 时，\(y-1=-\frac1{x-1}>0\)，故图象在直线 \(x=1\) 左侧且位于 \(y=1\) 上方；当 \(x>1\) 时，图象在 \(y=1\) 下方. 与选项图形比较可知应选 B.
\end{solution}

\begin{problem}
从正方体的 \(6\) 个面中选取 \(3\) 个面，其中有 \(2\) 个面不相邻的选法共有（\quad）

\choices
    {\(8\) 种}
    {\(12\) 种}
    {\(16\) 种}
    {\(20\) 种}
\end{problem}

\begin{answer}
B
\end{answer}

\begin{solution}
正方体的三个相对面共有 \(3\) 对. 先选定一对不相邻的相对面，再从其余 \(4\) 个面中选取第三个面，共有 \(3\times4=12\) 种. 一个三面组合不可能同时含有两对相对面，所以没有重复计数，故选 B.
\end{solution}

\begin{problem}
据 \(2002\text{ 年}\) \(3\text{ 月}\) \(5\text{ 日}\)九届人大五次会议《政府工作报告》：“\(2001\text{ 年}\)国内生产总值达到 \(95933\text{ 亿元}\)，比上年增长 \(7.3\%\)”，如果“十·五”期间（\(2001\text{ 年}\)－\(2005\text{ 年}\)）每年的国内生产总值都按此年增长率增长，那么到“十·五”末我国国内年生产总值约为（\quad）

\choices
    {\(115000\text{ 亿元}\)}
    {\(120000\text{ 亿元}\)}
    {\(127000\text{ 亿元}\)}
    {\(135000\text{ 亿元}\)}
\end{problem}

\begin{answer}
C
\end{answer}

\begin{solution}
以 \(2001\) 年末国内生产总值为基数，经过 \(2002\) 年至 \(2005\) 年共四次增长后，\(2005\) 年末的国内生产总值为
\(\displaystyle 95933(1+7.3\%)^4=95933(1.073)^4\). 又 \(1.073^2=1.151329\)，所以 \(1.073^4\approx1.32556\)，从而 \(95933(1.073)^4\approx127165\)（亿元），约为 \(127000\) 亿元，故选 C.
\end{solution}

\section{填空题}
\begin{problem}
函数 \( y = a^x \) 在 \([0,1]\) 上的最大值与最小值之和为 3，则 \( a = \)\fillinblank{}.
\end{problem}

\begin{answer}
2
\end{answer}

\begin{solution}
指数函数的底数满足 \(a>0\) 且 \(a\neq1\). 当 \(a>1\) 时，\(y=a^x\) 在 \([0,1]\) 上递增，最小值为 \(a^0=1\)，最大值为 \(a^1=a\)，故 \(a+1=3\)，得 \(a=2\). 当 \(0<a<1\) 时，函数在 \([0,1]\) 上递减，最大值与最小值之和为 \(1+a<2\)，不可能等于 \(3\). 因此 \(a=2\).
\end{solution}

\begin{problem}
椭圆 \(5x^{2} + ky^{2} = 5\) 的一个焦点是 \((0,2)\)，那么 \(k = \)\fillinblank{}.
\end{problem}

\begin{answer}
1
\end{answer}

\begin{solution}
将椭圆方程化为 \(\displaystyle\frac{x^2}{1}+\frac{y^2}{5/k}=1\). 焦点在 \(y\) 轴上，故长半轴平方为 \(\displaystyle a^2=\frac{5}{k}\)，短半轴平方为 \(b^2=1\). 由焦点为 \((0,2)\)，有 \(c=2\)，从而 \(c^2=a^2-b^2\)，即 \(4=\frac{5}{k}-1\). 解得 \(k=1\).
\end{solution}

\begin{problem}
\((x^{2} + 1)(x - 2)^{7}\) 展开式中 \(x^{3}\) 的系数是\fillinblank{}.
\end{problem}

\begin{answer}
1008
\end{answer}

\begin{solution}
要得到 \(x^3\) 项，只有两种取法：从 \(x^2(x-2)^7\) 中取 \((x-2)^7\) 的一次项，贡献 \(\displaystyle x^2\cdot\mathrm C_7^1x(-2)^6\)；从 \((x-2)^7\) 中直接取三次项. 因此系数为
\(\displaystyle\mathrm C_7^1(-2)^6+\mathrm C_7^3(-2)^4=7\times64+35\times16=1008\).
\end{solution}

\begin{problem}
已知 \( f(x)=\frac{x^{2}}{1+x^{2}} \)，那么 \( f(1)+f(2)+f\left(\frac{1}{2}\right)+f(3)+f\left(\frac{1}{3}\right)+f(4)+f\left(\frac{1}{4}\right)= \)\fillinblank{}.
\end{problem}

\begin{answer}
\(\frac{7}{2}\)
\end{answer}

\begin{solution}
对任意非零 \(x\)，有
\(\displaystyle f(x)+f\left(\frac1x\right)=\frac{x^2}{1+x^2}+\frac{1/x^2}{1+1/x^2}=\frac{x^2}{1+x^2}+\frac1{1+x^2}=1\). 所以 \(f(2)+f\left(\frac12\right)=1\)，\(f(3)+f\left(\frac13\right)=1\)，\(f(4)+f\left(\frac14\right)=1\)，且 \(f(1)=\frac12\). 原式为 \(1+1+1+\frac12=\frac72\).
\end{solution}

\section{解答题}
\begin{problem}
已知 \(\sin^2 2\alpha + \sin 2\alpha \cos \alpha - \cos 2\alpha = 1\)（\(\alpha \in \left(0, \frac{\pi}{2}\right)\)），求 \(\sin \alpha\)，\(\tan \alpha\) 的值.
\end{problem}

\begin{solution}
令 \(s=\sin\alpha\)，\(c=\cos\alpha\). 因为 \(0<\alpha<\frac{\pi}{2}\)，所以 \(0<s<1\) 且 \(c>0\). 利用 \(\sin2\alpha=2sc\)、\(\cos2\alpha=c^2-s^2\) 和 \(c^2=1-s^2\)，原等式化为
\(\displaystyle4s^2(1-s^2)+2s(1-s^2)-(1-2s^2)=1\)，即 \(\displaystyle2s^4+s^3-3s^2-s+1=0\). 因式分解得 \(\displaystyle(2s-1)(s+1)^2(s-1)=0\). 结合 \(0<s<1\)，得 \(s=\frac12\)，从而 \(c=\frac{\sqrt3}{2}\)，所以 \(\sin\alpha=\frac12\)，\(\tan\alpha=\frac{s}{c}=\frac{\sqrt3}{3}\).
\end{solution}

\begin{problem}
正方形 \(ABCD\)，\(ABEF\) 的边长都是 \(1\)，而且平面 \(ABCD\)，\(ABEF\) 互相垂直. 点 \(M\) 在 \(AC\) 上移动，点 \(N\) 在 \(BF\) 上移动，若 \(CM = BN = a\)（\(0 < a < \sqrt{2}\)）.

\begin{enumerate}
\item 求 \(MN\) 的长；
\item \(a\) 为何值时，\(MN\) 的长最小；
\item 当 \(MN\) 的长最小时，求面 \(MNA\) 与面 \(MNB\) 所成二面角 \(\alpha\) 的大小.
\end{enumerate}
\bitmapfigure[max width=0.36\linewidth]{img/2002/syllabus_science/q18_fig1.png}
\end{problem}

\begin{solution}
\begin{enumerate}
\item 取 \(A\) 为原点，令 \(AB\)、\(AD\)、\(AF\) 分别为 \(x\)、\(y\)、\(z\) 轴正向，则 \(B(1,0,0)\)、\(C(1,1,0)\)、\(D(0,1,0)\)、\(F(0,0,1)\)、\(E(1,0,1)\). 因为 \(AC=BF=\sqrt2\)，且 \(CM=BN=a\)，所以 \(\displaystyle M\left(1-\frac{a}{\sqrt2},1-\frac{a}{\sqrt2},0\right)\)，\(\displaystyle N\left(1-\frac{a}{\sqrt2},0,\frac{a}{\sqrt2}\right)\). 因而
\(\displaystyle MN^2=\left(1-\frac{a}{\sqrt2}\right)^2+\left(\frac{a}{\sqrt2}\right)^2=a^2-\sqrt2a+1=\left(a-\frac{\sqrt2}{2}\right)^2+\frac12\)，故 \(\displaystyle MN=\sqrt{\left(a-\frac{\sqrt2}{2}\right)^2+\frac12}\).
\item 由上式，\(\displaystyle\left(a-\frac{\sqrt2}{2}\right)^2\geqslant0\)，所以当 \(\displaystyle a=\frac{\sqrt2}{2}\) 时，\(MN\) 取得最小值 \(\displaystyle\frac{\sqrt2}{2}\).
\item 当 \(\displaystyle a=\frac{\sqrt2}{2}\) 时，\(M\left(\frac12,\frac12,0\right)\)，\(N\left(\frac12,0,\frac12\right)\). 设 \(H\) 为 \(MN\) 的中点，则 \(H\left(\frac12,\frac14,\frac14\right)\). 由坐标直接计算，\(\overrightarrow{MN}=(0,-\frac12,\frac12)\)、\(\overrightarrow{HA}=(-\frac12,-\frac14,-\frac14)\)、\(\overrightarrow{HB}=(\frac12,-\frac14,-\frac14)\)，于是 \(\overrightarrow{HA}\cdot\overrightarrow{MN}=\overrightarrow{HB}\cdot\overrightarrow{MN}=0\). 因而 \(HA\perp MN\)、\(HB\perp MN\)，且 \(\angle AHB\) 是所求二面角的平面角. 又 \(\displaystyle HA^2=HB^2=\frac38\)，\(AB=1\)，由余弦定理，
\(\displaystyle\cos\alpha=\cos\angle AHB=\frac{HA^2+HB^2-AB^2}{2HA\cdot HB}=\frac{\frac38+\frac38-1}{2\cdot\frac38}=-\frac13\). 因此 \(\displaystyle\alpha=\arccos\left(-\frac13\right)\).
\end{enumerate}
\end{solution}

\begin{problem}
设点 \(P\) 到点 \((-1,0)\)，\((1,0)\) 距离之差为 \(2m\)，到 \(x\) 轴、\(y\) 轴的距离之比为 2，求 \(m\) 的取值范围.
\end{problem}

\begin{solution}
设 \(P(x,y)\)，则到 \(x\) 轴、\(y\) 轴的距离分别为 \(|y|\)、\(|x|\). 由题意，\(\displaystyle\frac{|y|}{|x|}=2\)，所以 \(y=\pm2x\)，且 \(x\neq0\). 于是 \(P\) 不在两焦点 \((-1,0)\)、\((1,0)\) 的连线上，故两焦点到 \(P\) 的距离之差的绝对值严格小于焦点间距离 \(2\)，从而 \(0<|m|<1\).

记 \(d_1=\sqrt{(x+1)^2+y^2}\)、\(d_2=\sqrt{(x-1)^2+y^2}\). 由题意 \(|d_1-d_2|=2|m|\)，所以以 \((-1,0)\)、\((1,0)\) 为焦点的双曲线的实半轴为 \(|m|\)，焦半距为 \(1\)，虚半轴平方为 \(1-m^2\). 因而 \(P\) 满足双曲线方程 \(\displaystyle\frac{x^2}{m^2}-\frac{y^2}{1-m^2}=1\). 代入 \(y^2=4x^2\)，得
\(\displaystyle x^2\left(\frac1{m^2}-\frac4{1-m^2}\right)=1\)，即 \(\displaystyle x^2=\frac{m^2(1-m^2)}{1-5m^2}\). 因为 \(x^2>0\)、\(0<m^2<1\)，所以必须有 \(1-5m^2>0\)，即 \(0<m^2<\frac15\).

反之，若 \(0<m^2<\frac15\)，则上式给出正数 \(x^2\)，取任意满足该式的非零 \(x\)，再取 \(y=2x\) 或 \(y=-2x\)，所得点满足距离比条件和双曲线方程，因而满足原距离差条件. 所以 \(m\) 的取值范围为 \(\displaystyle\left(-\frac{\sqrt5}{5},0\right)\cup\left(0,\frac{\sqrt5}{5}\right)\).
\end{solution}

\begin{problem}
某城市 \(2001\text{ 年}\)末汽车保有量为 \(30\text{ 万辆}\)，预计此后每年报废上一年末汽车保有量的 \(6\%\)，并且每年新增汽车数量相同. 为保护城市环境，要求该城市汽车保有量不超过 \(60\text{ 万辆}\)，那么每年新增汽车数量不应超过多少辆？
\end{problem}

\begin{solution}
以万辆为单位，设每年新增汽车数量为 \(x\)（\(x\geqslant0\)），记 \(u_n\) 为第 \(n\) 年末汽车保有量，其中 \(u_0=30\). 由每年报废上一年末保有量的 \(6\%\)，得 \(\displaystyle u_{n+1}=0.94u_n+x\). 令 \(\displaystyle L=\frac{x}{1-0.94}=\frac{x}{0.06}\)，则
\(\displaystyle u_{n+1}-L=0.94(u_n-L)\)，所以 \(\displaystyle u_n=L+(30-L)0.94^n\).

当 \(0\leqslant x\leqslant1.8\) 时，\(L\leqslant30\)，于是 \(\displaystyle u_n=L+(30-L)0.94^n\leqslant L+(30-L)=30\). 当 \(1.8<x\leqslant3.6\) 时，\(30<L\leqslant60\)，且 \(30-L<0\)，所以 \(u_n<L\leqslant60\)，并且保有量随 \(n\) 增加而递增，始终不超过 \(60\).

反之，若 \(x>3.6\)，则 \(L>60\)，且 \(u_n=L+(30-L)0.94^n\) 递增并趋于 \(L>60\). 因为数列收敛到大于 \(60\) 的数，所以存在某个 \(n\) 使 \(u_n>60\)，不能满足题设要求. 因此 \(x\leqslant3.6\) 万辆，即每年新增汽车数量不应超过 \(36000\) 辆.
\end{solution}

\begin{problem}
设 \(a\) 为实数，函数 \(f(x) = x^{2} + |x - a| + 1\)（\(x \in \R\)）.

\begin{enumerate}
\item 讨论 \(f(x)\) 的奇偶性；
\item 求 \(f(x)\) 的最小值.
\end{enumerate}
\end{problem}

\begin{solution}
\begin{enumerate}
\item 因为 \(f(0)=|a|+1>0\)，所以 \(f(x)\) 不可能是奇函数. 当 \(a=0\) 时，\(f(-x)=x^2+|-x|+1=f(x)\)，故为偶函数. 若 \(a\neq0\)，取 \(x=a\)，则 \(f(a)=a^2+1\)，而 \(f(-a)=a^2+2|a|+1\neq f(a)\)，故不是偶函数. 所以当且仅当 \(a=0\) 时为偶函数，其他情形既非奇函数也非偶函数.
\item 按 \(x\) 与 \(a\) 的大小去掉绝对值符号：当 \(x\leqslant a\) 时，
\(\displaystyle f(x)=x^2-x+a+1=\left(x-\frac12\right)^2+a+\frac34\)；当 \(x\geqslant a\) 时，
\(\displaystyle f(x)=x^2+x-a+1=\left(x+\frac12\right)^2-a+\frac34\).

在区间 \((-\infty,a]\) 上，第一段平方项的顶点为 \(x=\frac12\)：若 \(a\leqslant\frac12\)，该段最小值在 \(x=a\) 处取得，为 \(a^2+1\)；若 \(a>\frac12\)，该段最小值在 \(x=\frac12\) 处取得，为 \(a+\frac34\). 在区间 \([a,+\infty)\) 上，第二段平方项的顶点为 \(x=-\frac12\)：若 \(a\leqslant-\frac12\)，该段最小值在 \(x=-\frac12\) 处取得，为 \(\frac34-a\)；若 \(a>-\frac12\)，该段最小值在 \(x=a\) 处取得，为 \(a^2+1\).

当 \(a\leqslant-\frac12\) 时，\(\displaystyle a^2+1-\left(\frac34-a\right)=\left(a+\frac12\right)^2\geqslant0\)，故全局最小值为 \(\frac34-a\). 当 \(-\frac12<a\leqslant\frac12\) 时，两段的最小值都在连接点 \(x=a\) 处取得，故全局最小值为 \(a^2+1\). 当 \(a>\frac12\) 时，\(\displaystyle a^2+1-\left(a+\frac34\right)=\left(a-\frac12\right)^2\geqslant0\)，故全局最小值为 \(a+\frac34\). 因此
\(\displaystyle\min_{x\in\R}f(x)=\begin{cases}
\frac34-a,&a\leqslant-\frac12,\\
a^2+1,&-\frac12<a\leqslant\frac12,\\
a+\frac34,&a>\frac12.
\end{cases}\)
\end{enumerate}
\end{solution}

\begin{problem}
设数列 \(\{a_{n}\}\) 满足： \(a_{n + 1} = a_n^2 - n a_n + 1\)（\(n = 1\)，\(2\)，\(3\)，\(\dots\)）.

\begin{enumerate}
\item 当 \(a_{1} = 2\) 时，求 \(a_{2}\)，\(a_{3}\)，\(a_{4}\) 并由此猜测 \(a_{n}\) 的一个通项公式；
\item 当 \(a_{1} \geqslant 3\) 时，证明对所有的 \(n \geqslant 1\)，有
\begin{circlelist}
\item \(a_{n}\geqslant n + 2\)；
\item \(\frac{1}{1 + a_1} +\frac{1}{1 + a_2} +\frac{1}{1 + a_3} +\dots +\frac{1}{1 + a_n}\leqslant \frac{1}{2}\).
\end{circlelist}
\end{enumerate}
\end{problem}

\begin{solution}
\begin{enumerate}
\item 由递推式，\(\displaystyle a_2=2^2-1\cdot2+1=3\)，\(\displaystyle a_3=3^2-2\cdot3+1=4\)，\(\displaystyle a_4=4^2-3\cdot4+1=5\). 由此猜测 \(a_n=n+1\). 下面用数学归纳法验证：当 \(n=1\) 时，\(a_1=2=1+1\)；若 \(a_n=n+1\)，则由递推式
\(\displaystyle a_{n+1}=a_n^2-na_n+1=(n+1)^2-n(n+1)+1=n+2=(n+1)+1\). 因而对所有 \(n\geqslant1\)，\(a_n=n+1\).
\item 先证明 \(a_n\geqslant n+2\). 当 \(n=1\) 时，题设给出 \(a_1\geqslant3=1+2\). 若 \(a_n\geqslant n+2\)，则 \(a_n-n\geqslant2\)，从而
\(\displaystyle a_{n+1}=a_n(a_n-n)+1\geqslant2a_n+1\geqslant2(n+2)+1\geqslant n+3\). 由数学归纳法，\(a_n\geqslant n+2\) 对所有 \(n\geqslant1\) 成立.

由上面的不等式，对每个 \(n\geqslant1\) 都有 \(a_n-n\geqslant2\)，所以 \(\displaystyle a_{n+1}=a_n(a_n-n)+1\geqslant2a_n+1\). 因而
\(\displaystyle 1+a_{n+1}\geqslant2(1+a_n)\)，即 \(\displaystyle\frac1{1+a_{n+1}}\leqslant\frac1{2(1+a_n)}\). 逐次迭代得到，对每个 \(k\geqslant1\)，
\(\displaystyle\frac1{1+a_k}\leqslant\frac1{2^{k-1}(1+a_1)}\). 于是
\(\displaystyle\sum_{k=1}^{n}\frac1{1+a_k}\leqslant\frac1{1+a_1}\sum_{k=1}^{n}\frac1{2^{k-1}}=\frac2{1+a_1}\left(1-\frac1{2^n}\right)\leqslant\frac2{1+a_1}\leqslant\frac12\)，其中最后一个不等式用了 \(a_1\geqslant3\).
\end{enumerate}
\end{solution}
